\scanpage{101}%
\mspage{\handwritten{5}}%
al structure of a\deleted{\unsure{n}} would-be $\infty$-groupoid, endowed only with its ``structure\handwritten{ of} order $\leq n$\handwritten{»}. The idea is to give a\deleted{\unsure{n}} direct inductive construction of this sequence, by describing $C_0$, and the passage from $C_n$ to $C_{n+1}$ ($n\geq0$), namely from $n$-ary to $(n+1)$-ary structure. As for the meaning of ``universal structure'', \deleted{\unsure{it depends}} once a given structure species is at hand, it depends on the type of categories (described by the exactness properties one is assuming for these) one wants to take as carriers for the considered structure, and the type of exactness properties one assumes for the functor one allows between these. The choice depends partly on the particular species;\deleted{\unsure{at|it}} if it is an algebraic structure which can be described say by a handful of composition laws between a bunch of base sets (or base-objects, when looking at ``realizations'' of the structure not only in the category (Sets)), one natural choice is to take categories with finite products, and functors which commute to these. For more sophisticated algebraic structures (including structure of category, groupoid or the like), which requires for the description of data or axioms not only finite products, but also some fiber products, one other familiar choice is to take categories with finite inverse limits, and left exact functors. \handwritten{Still} more sophisticated \deleted{choices}\inserted{\handwritten{structures}}, when the description of the structure in terms of base objects requires not only some kinds of inverse limits, but also more or less arbitrary direct limits (such as the structure of a comodule over an algebra, which requires consideration of tensor products over a ring object),%
% slanted
\margin{\handwritten{or the structure \unsure{of} \unsure{-|a} torsor under \unsure{a} group}} still \deleted{\unsure{more|much}} more stringent conditions will have to be imposed upon categories and functors between these, for the structure to make a sense in \deleted{any of} these categories, and the functors to transform a structure of this type in one category into one of same type in another. In most examples I have looked up, everything is \unsure{OK} taking categories which are topoi, and functors between these which are inverse image functors for morphisms of topoi, namely which are left exact and commute with arbitrary direct limits. Th\deleted{\unsure{er|re}}ere is a general theorem for the existence of universal structures, covering all these cases - for instance there is a ``classifying topos'' for most algebr-geometric structures, whose cohomology say should be viewed as the ``classifying cohomology'' of the structure species considered. In the case we are interested in here, it is convenient however to work with the smallest categories $C_n$ feasable - which amounts to being as generous as possible for the categories \deleted{and functors} one is allowing as carriers for the structure of an $\infty$-groupoid, and for the functors between these which are expected to carry an $\infty$-groupoid into an $\infty$-groupoid. What we will do is define ultimately a structure of an $\infty$-groupoid in a category $C$, as a sequence of objects $F_i$ ($i\in\underline{N}$), endowed with some structure to be defined,
% archive sheet number: 100
